% Generated draft: requires mathematical and source review.
\documentclass[11pt]{article}
\usepackage[margin=25mm]{geometry}
\usepackage{fontspec}
\setmainfont{Latin Modern Roman}
\usepackage{amsmath,amssymb,mathtools,unicode-math}
\setmathfont{Latin Modern Math}
\usepackage{enumitem,needspace,xurl,hyperref}
\setlength{\parindent}{0pt}
\setlength{\parskip}{5pt}
\setlength{\emergencystretch}{3em}
\newcommand{\N}{\mathbb{N}}
\newcommand{\Z}{\mathbb{Z}}
\newcommand{\Q}{\mathbb{Q}}
\newcommand{\R}{\mathbb{R}}
\newcommand{\C}{\mathbb{C}}
\newcommand{\F}{\mathbb{F}}
\newcommand{\A}{\mathbb{A}}
\newcommand{\PP}{\mathbb{P}}
\newcommand{\E}{\mathbb{E}}
\newcommand{\eps}{\varepsilon}
\newcommand{\dd}{\,\mathrm{d}}
\begin{document}
\section*{UnsolvedMath ID: 30000080}
\textbf{Source problem code:} catalogue entry
\subsection{Definitions and mathematical statement}
Let $K$ be a field, $n\geq1$, and $S=K[x_1,\ldots,x_n]$ with its standard grading and homogeneous maximal ideal $\mathfrak m=(x_1,\ldots,x_n)$. Let $I\subsetneq S$ be a homogeneous $\mathfrak m$-primary ideal, meaning $\sqrt I=\mathfrak m$, generated in one degree $d\geq1$, so $I=S I_d$. Say that $I$ is linearly presented if it has a graded presentation
\[
 S(-d-1)^b\longrightarrow S(-d)^a\longrightarrow I\longrightarrow0
\]
whose first map has linear-form entries; equivalently, the first syzygies on minimal degree-$d$ generators are generated in degree $d+1$. This imposes no claim that every higher syzygy is linear.

The conjecture is
\[
 I^{\,n-1}=\mathfrak m^{(n-1)d}
\]
for every ideal satisfying these assumptions. Powers are ordinary ideal powers. Equivalently, since $I^{n-1}$ is generated in degree $(n-1)d$, all forms of that degree must lie in it:
\[
 (I^{n-1})_{(n-1)d}=S_{(n-1)d}.
\]
For $n=1$ the assertion is the convention $I^0=S=\mathfrak m^0$; the substantive exponent is positive for $n\geq2$. No characteristic restriction on $K$ is supplied.

The supplied literature assessment checked 21 August 2026 reports the precise $(n-1)$-power equality proved for $n=3$ and for monomial ideals, while describing the general result only as eventual equality: for a given ideal, $I^t=\mathfrak m^{td}$ for all sufficiently large $t$. Such an ideal-dependent eventual threshold is weaker than the uniform threshold $t=n-1$. The assessment calls the all-dimensional conjecture open. These are dated supplied claims; full theorem hypotheses, including any field qualifications in the references, require review.
\subsection{Short English statement}
Must the $(n-1)$st ordinary power of a degree-$d$, linearly presented primary ideal in $n$ variables contain every form of degree $(n-1)d$? The supplied assessment reports the bound in dimension three and for monomial ideals, but only eventual equality in general.
\subsection{Sources}
\begin{itemize}
\item Local UnsolvedMath registry, permanent ID 30000080, exact source code catalogue entry; \emph{Powers of Linearly Presented Primary Ideals}. Proposed year 2004, status open; corrected\_verified extraction retains the displayed conjecture.
\item Supplied report citation: \emph{Combinatorial Commutative Algebra (2005)}, Oberwolfach Reports, DOI 10.4171/owr/2004/32, pp. 1703--1746, \url{https://doi.org/10.4171/owr/2004/32}. The title year differs from DOI/proposed year.
\item Supplied literature assessment checked 2026-08-21: Eisenbud--Huneke--Ulrich, \emph{The Regularity of Tor}, \url{https://arxiv.org/abs/math/0405373}; ``Published version'', \url{https://doi.org/10.1353/ajm.2006.0022}; ``Original Oberwolfach report'', \url{https://doi.org/10.4171/owr/2004/32}. This dated OWR-derived triage requests independent verification of time-sensitive status and cited claims.
\end{itemize}
\end{document}
